\documentclass{amsart}
% For dealing with references we use the comment environment
\usepackage{verbatim}
\newenvironment{reference}{\comment}{\endcomment}
%\newenvironment{reference}{}{}
\newenvironment{slogan}{\comment}{\endcomment}
\newenvironment{history}{\comment}{\endcomment}

% For commutative diagrams we use Xy-pic
\usepackage[all]{xy}

% We use 2cell for 2-commutative diagrams.
\xyoption{2cell}
\UseAllTwocells

% We use multicol for the list of chapters between chapters
\usepackage{multicol}

% This is generally recommended for better output
\usepackage{lmodern}
\usepackage[T1]{fontenc}

% For cross-file-references

% Package for hypertext links:
\usepackage{hyperref}

% For any local file, say "hello.tex" you want to link to please

% Theorem environments.
%
\theoremstyle{plain}
\newtheorem{theorem}[subsection]{Theorem}
\newtheorem{proposition}[subsection]{Proposition}
\newtheorem{lemma}[subsection]{Lemma}

\theoremstyle{definition}
\newtheorem{definition}[subsection]{Definition}
\newtheorem{example}[subsection]{Example}
\newtheorem{exercise}[subsection]{Exercise}
\newtheorem{situation}[subsection]{Situation}

\theoremstyle{remark}
\newtheorem{remark}[subsection]{Remark}
\newtheorem{remarks}[subsection]{Remarks}

\numberwithin{equation}{subsection}

% Macros
%
\def\lim{\mathop{\mathrm{lim}}\nolimits}
\def\colim{\mathop{\mathrm{colim}}\nolimits}
\def\Spec{\mathop{\mathrm{Spec}}}
\def\Hom{\mathop{\mathrm{Hom}}\nolimits}
\def\Ext{\mathop{\mathrm{Ext}}\nolimits}
\def\SheafHom{\mathop{\mathcal{H}\!\mathit{om}}\nolimits}
\def\SheafExt{\mathop{\mathcal{E}\!\mathit{xt}}\nolimits}
\def\Sch{\mathit{Sch}}
\def\Mor{\mathop{\mathrm{Mor}}\nolimits}
\def\Ob{\mathop{\mathrm{Ob}}\nolimits}
\def\Sh{\mathop{\mathit{Sh}}\nolimits}
\def\NL{\mathop{N\!L}\nolimits}
\def\CH{\mathop{\mathrm{CH}}\nolimits}
\def\proetale{{pro\text{-}\acute{e}tale}}
\def\etale{{\acute{e}tale}}
\def\QCoh{\mathit{QCoh}}
\def\Ker{\mathop{\mathrm{Ker}}}
\def\Im{\mathop{\mathrm{Im}}}
\def\Coker{\mathop{\mathrm{Coker}}}
\def\Coim{\mathop{\mathrm{Coim}}}

% Boxtimes
%
\DeclareMathSymbol{\boxtimes}{\mathbin}{AMSa}{"02}

%
% Macros for moduli stacks/spaces
%
\def\QCohstack{\mathcal{QC}\!\mathit{oh}}
\def\Cohstack{\mathcal{C}\!\mathit{oh}}
\def\Spacesstack{\mathcal{S}\!\mathit{paces}}
\def\Quotfunctor{\mathrm{Quot}}
\def\Hilbfunctor{\mathrm{Hilb}}
\def\Curvesstack{\mathcal{C}\!\mathit{urves}}
\def\Polarizedstack{\mathcal{P}\!\mathit{olarized}}
\def\Complexesstack{\mathcal{C}\!\mathit{omplexes}}
% \Pic is the operator that assigns to X its picard group, usage \Pic(X)
% \Picardstack_{X/B} denotes the Picard stack of X over B
% \Picardfunctor_{X/B} denotes the Picard functor of X over B
\def\Pic{\mathop{\mathrm{Pic}}\nolimits}
\def\Picardstack{\mathcal{P}\!\mathit{ic}}
\def\Picardfunctor{\mathrm{Pic}}
\def\Deformationcategory{\mathcal{D}\!\mathit{ef}}

\setlength{\emergencystretch}{3em}
\begin{document}
\title[Stacks Project, Tag 06DM]{1746 --- An arbitrary local ring is geometrically unibranch exactly when its strict henselization has a unique minimal prime}
\maketitle
\noindent Copyright (C) 2005--2025 Johan de Jong. Stacks Project authors. Source: \href{https://stacks.math.columbia.edu/tag/06DM}{Tag 06DM}.

\noindent Editorial difficulty note (not author text). Difficulty H1: The proof characterizes geometric unibranchness for arbitrary local rings by the strict henselization rather than assuming a finite normalization or Noetherian ring. It constructs finite normalization subalgebras witnessing either disconnected branches or a non-purely-inseparable residue extension, forces a forbidden product inside a domain after strict henselian base change, and proves the converse by normality/locality of every etale local stage. This normalization-versus-geometric-branch argument is above ordinary qualification algebra..

\noindent Editorial provenance note (not author text). History: excerpt from revision \texttt{a0444\allowbreak 6e57e\allowbreak c1fbc\allowbreak 252a8\allowbreak 71afc\allowbreak ec775\allowbreak 2fb28\allowbreak 07b14}, more-algebra.tex, lines 31452--31550. Reference presentation was adapted; original-author context and core support blocks were retained or added. The mathematical wording is unchanged. Licensed under GNU FDL 1.2 or later; no invariant sections or cover texts. See ../licenses/COPYING. Context blocks reproduce author definitions and supporting results; they are not additional counted problems.

\begin{definition}
\label{definition-unibranch}
\begin{reference}
\href{https://stacks.math.columbia.edu/bibliography}{Bibliography: EGA4 (Chapter 0 (23.2.1))}
\end{reference}
Let $A$ be a local ring. We say $A$ is {\it unibranch}
if the reduction $A_{red}$ is a domain and if the integral closure
$A'$ of $A_{red}$ in its field of fractions is local.
We say $A$ is {\it geometrically unibranch} if $A$ is unibranch
and moreover the residue field of $A'$ is purely inseparable over
the residue field of $A$.
\end{definition}

\begin{definition}
\label{algebra-definition-henselian}
Let $(R, \mathfrak m, \kappa)$ be a local ring.
\begin{enumerate}
\item We say $R$ is {\it henselian} if for every monic $f \in R[T]$ and
every root $a_0 \in \kappa$ of $\overline{f}$ such that
$\overline{f'}(a_0) \not = 0$
there exists an $a \in R$ such that $f(a) = 0$ and
$a_0 = \overline{a}$.
\item We say $R$ is {\it strictly henselian} if $R$ is henselian
and its residue field is separably algebraically closed.
\end{enumerate}
\end{definition}

\begin{definition}
\label{algebra-definition-ring-normal}
A ring $R$ is called {\it normal} if for every prime
$\mathfrak p \subset R$ the localization $R_{\mathfrak p}$ is
a normal domain (see Definition \ref{algebra-definition-domain-normal}).
\end{definition}

\begin{definition}
\label{algebra-definition-domain-normal}
A domain $R$ is called {\it normal} if it is integrally
closed in its field of fractions.
\end{definition}

\begin{lemma}
\label{lemma-henselization-reduced}
\begin{slogan}
Reducedness passes to the (strict) henselization.
\end{slogan}
Let $R$ be a local ring.
The following are equivalent: $R$ is reduced,
the henselization $R^h$ of $R$ is reduced, and
the strict henselization $R^{sh}$ of $R$ is reduced.
\end{lemma}

\begin{proof}
The ring maps $R \to R^h \to R^{sh}$ are faithfully flat.
Hence one direction of the implications follows from
Algebra, Lemma \href{https://stacks.math.columbia.edu/tag/033F}{lemma descent reduced (Tag 033F)}.
Conversely, assume $R$ is reduced. Since $R^h$ and $R^{sh}$
are filtered colimits of \'etale, hence smooth $R$-algebras, the
result follows from
Algebra, Lemma \href{https://stacks.math.columbia.edu/tag/033B}{lemma reduced goes up (Tag 033B)}.
\end{proof}


\begin{lemma}
\label{algebra-lemma-local-tensor-with-integral}
Let $A \to B$ and $A \to C$ be local homomorphisms of local rings.
If $A \to C$ is integral and either
$\kappa(\mathfrak m_C)/\kappa(\mathfrak m_A)$ or
$\kappa(\mathfrak m_B)/\kappa(\mathfrak m_A)$ is purely
inseparable, then $D = B \otimes_A C$ is a local ring and
$B \to D$ and $C \to D$ are local.
\end{lemma}

\begin{proof}
Any maximal ideal of $D$ lies over the maximal ideal of $B$ by
going up for the integral ring map
$B \to D$ (Lemma \href{https://stacks.math.columbia.edu/tag/00GU}{lemma integral going up (Tag 00GU)}).
Now $D/\mathfrak m_B D = \kappa(\mathfrak m_B) \otimes_A C =
\kappa(\mathfrak m_B) \otimes_{\kappa(\mathfrak m_A)} C/\mathfrak m_A C$.
The spectrum of $C/\mathfrak m_A C$ consists of a
single point, namely $\mathfrak m_C$. Thus the spectrum of
$D/\mathfrak m_B D$ is the same as the spectrum of
$\kappa(\mathfrak m_B) \otimes_{\kappa(\mathfrak m_A)} \kappa(\mathfrak m_C)$
which is a single point by our assumption that either
$\kappa(\mathfrak m_C)/\kappa(\mathfrak m_A)$ or
$\kappa(\mathfrak m_B)/\kappa(\mathfrak m_A)$ is purely
inseparable. This proves that $D$ is local and that the ring maps
$B \to D$ and $C \to D$ are local.
\end{proof}


\begin{lemma}
\label{algebra-lemma-mop-up-strictly-henselian}
Let $(R, \mathfrak m, \kappa)$ be a strictly henselian local ring.
Any finite type $R$-algebra $S$ can be written as
$S = A_1 \times \ldots \times A_n \times B$ with $A_i$ local
and finite over $R$ and $\kappa \subset \kappa(\mathfrak m_{A_i})$
finite purely inseparable and $R \to B$ not quasi-finite
at any prime of $B$ lying over $\mathfrak m$.
\end{lemma}

\begin{proof}
First write $S = A_1 \times \ldots \times A_n \times B$ as in
Lemma \href{https://stacks.math.columbia.edu/tag/04GJ}{lemma mop up (Tag 04GJ)}.
The field extension $\kappa(\mathfrak m_{A_i})/\kappa$
is finite and $\kappa$ is separably algebraically closed, hence
it is finite purely inseparable.
\end{proof}


\noindent
Let $A$ be a local ring. Here is an equivalent formulation
\begin{enumerate}
\item $A$ is unibranch if $A$ has a unique minimal prime $\mathfrak p$ and
the integral closure of $A/\mathfrak p$ in its fraction field is a
local ring, and
\item $A$ is geometrically unibranch if $A$ has a unique minimal prime
$\mathfrak p$ and the integral closure of $A/\mathfrak p$ in its
fraction field is a local ring whose residue field is purely inseparable
over the residue field of $A$.
\end{enumerate}
A local ring which is normal is geometrically unibranch
(follows from Definition \ref{definition-unibranch} and
Algebra, Definition \ref{algebra-definition-ring-normal}).
Lemmas \href{https://stacks.math.columbia.edu/tag/0BQ0}{lemma unibranch (Tag 0BQ0)} and \ref{lemma-geometrically-unibranch}
suggest that being (geometrically) unibranch
is a reasonable property to look at.


\begin{lemma}
\label{lemma-geometrically-unibranch}
\begin{reference}
\href{https://stacks.math.columbia.edu/bibliography}{Bibliography: Etale-coverings (Lemma 2.2)} and
\href{https://stacks.math.columbia.edu/bibliography}{Bibliography: EGA4 (Chapter IV Proposition 18.8.15)}
\end{reference}
Let $A$ be a local ring. Let $A^{sh}$ be a strict henselization of $A$.
The following are equivalent
\begin{enumerate}
\item $A$ is geometrically unibranch, and
\item $A^{sh}$ has a unique minimal prime.
\end{enumerate}
\end{lemma}

\begin{proof}
This follows from Lemma \href{https://stacks.math.columbia.edu/tag/0C25}{lemma geometric branches (Tag 0C25)}
but we will also give a direct proof; this direct
proof is almost exactly the same as the direct proof of
Lemma \href{https://stacks.math.columbia.edu/tag/0BQ0}{lemma unibranch (Tag 0BQ0)}.
Denote $\mathfrak m$ the maximal ideal of the ring $A$.
Denote $\kappa$, $\kappa^{sh}$ the residue field of $A$, $A^{sh}$.

\medskip\noindent
Assume (2). Let $\mathfrak p^{sh}$ be the unique minimal prime of
$A^{sh}$. The flatness of $A \to A^{sh}$ implies that
$\mathfrak p = A \cap \mathfrak p^{sh}$ is the unique minimal
prime of $A$ (by going down, see
Algebra, Lemma \href{https://stacks.math.columbia.edu/tag/00HS}{lemma flat going down (Tag 00HS)}).
Also, since $A^{sh}/\mathfrak pA^{sh} = (A/\mathfrak p)^{sh}$ (see
Algebra, Lemma \href{https://stacks.math.columbia.edu/tag/05WS}{lemma quotient strict henselization (Tag 05WS)})
is reduced by Lemma \ref{lemma-henselization-reduced}
we see that $\mathfrak p^{sh} = \mathfrak pA^{sh}$.
Let $A'$ be the integral closure of $A/\mathfrak p$ in its fraction
field. We have to show that $A'$ is local and that its residue
field is purely inseparable over $\kappa$.
Since $A \to A'$ is integral, every maximal ideal of $A'$ lies over
$\mathfrak m$ (by going up for integral ring maps, see
Algebra, Lemma \href{https://stacks.math.columbia.edu/tag/00GU}{lemma integral going up (Tag 00GU)}).
If $A'$ is not local, then we can find distinct maximal ideals
$\mathfrak m_1$, $\mathfrak m_2$. Choosing elements $f_1, f_2 \in A'$
with $f_i \in \mathfrak m_i, f_i \not \in \mathfrak m_{3 - i}$ we find
a finite subalgebra $B = A[f_1, f_2] \subset A'$ with distinct maximal
ideals $B \cap \mathfrak m_i$, $i = 1, 2$. If $A'$ is local with maximal
ideal $\mathfrak m'$, but $A/\mathfrak m \subset A'/\mathfrak m'$
is not purely inseparable, then we can find $f \in A'$ whose image in
$A'/\mathfrak m'$ generates a finite, not purely inseparable extension
of $A/\mathfrak m$ and we find a finite local subalgebra $B = A[f] \subset A'$
whose residue field is not a purely inseparable extension of $A/\mathfrak m$.
Note that the inclusions
$$
A/\mathfrak p \subset B \subset \kappa(\mathfrak p)
$$
give, on tensoring with the flat ring map $A \to A^{sh}$ the inclusions
$$
A^{sh}/\mathfrak p^{sh} \subset
B \otimes_A A^{sh} \subset
\kappa(\mathfrak p) \otimes_A A^{sh} \subset
\kappa(\mathfrak p^{sh})
$$
the last inclusion because
$\kappa(\mathfrak p) \otimes_A A^{sh} =
\kappa(\mathfrak p) \otimes_{A/\mathfrak p} A^{sh}/\mathfrak p^{sh}$
is a localization of the domain $A^{sh}/\mathfrak p^{sh}$.
Note that $B \otimes_A \kappa^{sh}$ has at least two maximal ideals
because $B/\mathfrak mB$ either has two maximal ideals or one whose
residue field is not purely inseparable over $\kappa$, and because
$\kappa^{sh}$ is separably algebraically closed. Hence, as
$A^{sh}$ is strictly henselian we see that
$B \otimes_A A^{sh}$ is a product of $\geq 2$ local rings, see
Algebra, Lemma \ref{algebra-lemma-mop-up-strictly-henselian}.
But we've just seen that $B \otimes_A A^{sh}$ is a subring of a domain
and we get a contradiction.

\medskip\noindent
Assume (1). Let $\mathfrak p \subset A$ be the unique minimal
prime and let $A'$ be the integral closure of $A/\mathfrak p$
in its fraction field. Let $A \to B$ be a local map of local rings which is a
localization of an \'etale $A$-algebra. In particular $\mathfrak m_B$
is the unique prime containing $\mathfrak m_AB$. Then $B' = A' \otimes_A B$
is integral over $B$ and the assumption that $A \to A'$ is local
with purely inseparable residue field extension implies that $B'$
is local (Algebra, Lemma \ref{algebra-lemma-local-tensor-with-integral}).
On the other hand, $A' \to B'$ is the localization
of an \'etale ring map, hence $B'$ is normal, see
Algebra, Lemma \href{https://stacks.math.columbia.edu/tag/033C}{lemma normal goes up (Tag 033C)}.
Thus $B'$ is a (local) normal domain. Finally, we have
$$
B/\mathfrak pB \subset B \otimes_A \kappa(\mathfrak p)
= B' \otimes_{A'} (\text{fraction field of }A') \subset
\text{fraction field of }B'
$$
Hence $B/\mathfrak pB$ is a domain, which implies that $B$ has a unique
minimal prime (since by flatness of $A \to B$ these all have to lie
over $\mathfrak p$). Since $A^{sh}$ is a filtered colimit of
the local rings $B$ it follows that $A^{sh}$ has a unique minimal prime.
Namely, if $fg = 0$ in $A^{sh}$ for some non-nilpotent elements
$f, g$, then we can find a $B$ as above containing both $f$ and $g$
which leads to a contradiction.
\end{proof}
\end{document}
